\documentclass[11pt]{article}
\usepackage{preamble}

\newcommand{\lessonrow}[2]{\textbf{#1} & #2 \\ \hline}

\newcommand{\placeholdertext}[1]{%
  {\small\itshape #1\par}
}

\newcommand{\chemistfigure}{%
\begin{center}
\begin{tikzpicture}[font=\small, line width=0.9pt]
\node[
  draw,
  rounded corners,
  fill=frameshade,
  minimum width=2.0in,
  minimum height=1.45in,
  align=center
] (chem) at (0,0) {\textbf{Chemist}\\pouring liquid\\into a container};
\node[
  draw,
  rounded corners,
  fill=calloutyellow,
  minimum width=2.0in,
  minimum height=0.62in,
  align=center
] (six) at (-2.45,-1.45) {\(x\) gallons of\\6\% alcohol};
\node[
  draw,
  rounded corners,
  fill=calloutpeach,
  minimum width=2.1in,
  minimum height=0.62in,
  align=center
] (two) at (2.45,-1.45) {50 gallons of\\2\% alcohol};
\draw[->, tealheader] (six.north east) .. controls (-1.6,-0.7) and (-0.8,-0.35) .. (chem.south west);
\draw[->, tealheader] (two.north west) .. controls (1.6,-0.7) and (0.8,-0.35) .. (chem.south east);
\end{tikzpicture}
\end{center}
}

\begin{document}

\packetheading{Lesson 4-5 $\cdot$ Quick $\cdot$ Student Packet}{Solving Rational Equations + Chemistry DOK-3 (Q\#2/Q\#6/Q\#7 Prep)}

\sectionbanner{OBJECTIVES}
{\small
\renewcommand{\arraystretch}{1.25}
\noindent\begin{tabular}{|p{1.8in}|p{4.8in}|}
\hline
\lessonrow{Math Objective}{Solve rational equations in one variable by clearing denominators, check for extraneous solutions, and apply rational equations to real-world rate and mixture problems.}
\lessonrow{Language Objective}{Justify why a candidate solution is extraneous, using the domain restrictions of the original rational equation.}
\lessonrow{Essential Understanding}{Multiplying both sides of a rational equation by an expression containing variables can introduce extraneous solutions --- values that satisfy the cleared equation but violate the original domain.}
\end{tabular}
}

\sectionbanner{TOPIC GOALS / ESSENTIAL QUESTION / MATERIALS}
{\small
\renewcommand{\arraystretch}{1.25}
\noindent\begin{tabular}{|p{1.8in}|p{4.8in}|}
\hline
\lessonrow{Topic Goals}{Fluency solving rational equations; identifying extraneous solutions; applying to work-rate and mixture problems (Topic 4 LEHS Q\#2, Q\#6, Q\#7 prep).}
\lessonrow{Essential Question}{When you solve a rational equation by clearing denominators, why must you check every candidate solution against the original equation?}
\lessonrow{Materials}{Student packet, pencil, calculator. DOK-3 Practice \#33 (chemistry mixture performance task) rehearses Topic 4 LEHS Q\#6 (solve rational equation) and Q\#7 (work-rate reasoning).}
\end{tabular}
}

\sectionbanner{LAUNCH --- Solving Rational Equations (teacher models Ex 1 + Ex 3)}

\begin{callout}{\IconBook\ RULES --- SOLVING RATIONAL EQUATIONS (keep printed on your page)}{calloutgreen}
1. To solve a rational equation:\\
\hspace*{1.6em}(a) find the LCD of all fractions,\\
\hspace*{1.6em}(b) multiply BOTH SIDES by the LCD to clear fractions,\\
\hspace*{1.6em}(c) solve the resulting polynomial equation,\\
\hspace*{1.6em}(d) CHECK each candidate in the ORIGINAL equation --- any candidate that makes a denominator zero is extraneous.\\
2. Domain restrictions are set BEFORE clearing denominators. A candidate that violates a restriction is extraneous even if it algebraically satisfies the cleared equation.\\
3. Work-rate setup: if person A does a job in a hours and person B in b hours, working together they complete \(1/a + 1/b\) of the job per hour, so combined time \(t\) satisfies \(1/a + 1/b = 1/t\).\\
4. Mixture setup: concentration \(\times\) volume = amount of substance. Build one fraction per mixture stream, then set up a rational equation for the target concentration.
\end{callout}

\begin{callout}{\IconWarn\ COMPRESSED LESSON --- TWO EXAMPLES IN LAUNCH}{calloutred}
This is a single-period quick treatment. Teacher models BOTH Example 1 (solve a rational equation) AND Example 3 (identify the extraneous solution) during Launch.\\
Pay attention to Example 3 --- the extraneous solution trap is the \#1 error on this lesson's assessment items. Always check every candidate in the ORIGINAL equation.
\end{callout}

\sectionbanner{EXPLORE --- Think-Pair-Share (32 min)}
\textit{Work with a partner. Move in order. Practice \#33 (Chemistry Mixture) is the big one --- plan \(\sim\)10 min for it.}

\textit{45-min Wed F variant: skip \#14 and \#25.}

\textbf{Bank-item labels (in order):}
\begin{enumerate}[leftmargin=1.6em,itemsep=2pt,topsep=2pt]
  \item Try It 1 --- Solve a rational equation
  \item Try It 3 --- Identify extraneous solution
  \item Practice \#10 --- Solve + check
  \item Practice \#14 --- Extraneous trap
  \item Practice \#25 --- Work-rate application
  \item Practice \#33 \IconStar\ DOK-3 CHEMISTRY MIXTURE --- Topic 4 LEHS Q\#6 + Q\#7 prep
\end{enumerate}

\bankitem{Try It 1 --- Solve a rational equation}{%
What is the solution to each equation?

\begin{enumerate}[label=\textbf{\alph*.}, leftmargin=1.6em, itemsep=3pt, topsep=3pt]
\item \(\frac{2}{x+5} = 4\)
\item \(\frac{1}{x-7} = 2\)
\end{enumerate}
}{}

\bankitem{Try It 3 --- Identify extraneous solution}{%
What is the solution to the equation \(\frac{1}{x+2} + \frac{1}{x-2} = \frac{4}{(x+2)(x-2)}\)?
}{}

\bankitem{Practice \#10 --- Solve + check}{%
Error Analysis Describe and correct the error Miranda made in solving the rational equation.
\[
\begin{aligned}
\frac{1}{x-2} + \frac{x-2}{x+2} &= \frac{x-4}{x-2} \\
(x+2)(x-2)\left(\frac{1}{x-2} + \frac{x-2}{x+2}\right) &= \left(\frac{x-4}{x-2}\right)(x+2)(x-2) \\
(x+2)(1) + (x-2)(x-2) &= (x+4)(x+2) \\
x + 2 + x^2 - 4x + 4 &= x^2 + 6x + 8 \\
-3x + 6 &= 6x + 8 \\
-2 &= 9x; \text{ or } x = -\frac{2}{9}\ \text{X}
\end{aligned}
\]
}{}

\bankitem{Practice \#14 --- Extraneous trap}{%
Make Sense and Persevere Solve the rational equation shown. Explain what is unique about the solution.
\[
\frac{x^2 - 7x - 18}{x+2} = x - 9
\]
}{}

\bankitem{Practice \#25 --- Work-rate application}{%
Make Sense and Persevere Kenji can finish a puzzle in 2 hours working alone. Oscar can finish the same puzzle in 3 hours working alone. How long would it take Oscar and Kenji to finish the puzzle if they worked on it together?
}{}

\bankitem{Practice \#33 \IconStar\ DOK-3 CHEMISTRY MIXTURE --- Topic 4 LEHS Q\#6 + Q\#7 prep}{%
\textbf{Performance Task}

A chemist needs alcohol solution in the correct concentration for her experiment. She adds a 6\% alcohol solution to 50 gallons of solution that is 2\% alcohol. The function that represents the percent of alcohol in the resulting solution is
\[
f(x) = \frac{50(0.02)+x(0.06)}{50+x},
\]
where \(x\) is the amount of 6\% solution added.

\chemistfigure
\placeholdertext{[IMAGE: Chemist pouring liquid into a container; labels ``x gallons of 6\% alcohol'' and ``50 gallons of 2\% alcohol.'']}

\textbf{Part A} How much 6\% solution should be added to create a solution that is 5\% alcohol?

\textbf{Part B} Use Appropriate Tools Explain the steps you could take to use your graphing calculator to verify the correctness of your answer to part (A).
}{}

\sectionbanner{SHARE / SUMMARY (5 min)}

\begin{sentenceframebox}
\textbf{Sentence frames}

Pair share (Chemistry): ``I set up \(f(x) = \frac{50\cdot0.02 + 0.06x}{50+x} = 0.05\). After clearing the denominator I got \_\_\_. Solving for \(x\) gave \(x =\) \_\_\_ gallons. I verified with a calculator TABLE by entering \(y_1 =\) \_\_\_ and \(y_2 =\) \_\_\_.''

Whole class: one pair reads their chemistry solution. Class signals thumbs up/sideways.
\end{sentenceframebox}

\begin{callout}{\IconSearch\ BRIDGE TO TOPIC 4 LEHS Q\#2, Q\#6, Q\#7}{calloutpeach}
The chemistry move --- set up rational equation, clear denominators, solve, check --- is exactly what Topic 4 LEHS Q\#6 asks. The domain-restriction reasoning is Q\#2. The rate reasoning is Q\#7. You have now rehearsed all three.
\end{callout}

\sectionbanner{EXIT}

\begin{summaryexitbox}{EXIT TICKET --- Summary Recap}
After clearing denominators in a rational equation, I must check each candidate solution in the \blank[1.1in] equation to catch \blank[1.0in] solutions.

In Practice \#33 (chemistry), the equation came from setting \blank[2.2in] equal to the target concentration 0.05.

One biggest thing I learned today: \dotfill
\end{summaryexitbox}

\clearpage

\sectionbanner{OPTIONAL REINFORCEMENT (not collected)}
\textit{If you finish Explore early. \#8 is a work-rate validity reasoning item that rehearses Topic 4 LEHS Q\#7 specifically.}

\bankitem{Optional \#1 (Savvas Practice \#8 --- work-rate validity reasoning)}{%
Reason If you solve a work-rate problem and your solution, which represents the amount of time it would take working together, exceeds the individual working alone times that are given, then how do you know your solution is unreasonable? Explain.
}{}

\end{document}
